\newpage
\begin{center}
  \textbf{\large 1. ЛИТЕРАТУРНЫЙ ОБЗОР}
\end{center}
\refstepcounter{chapter}
\addcontentsline{toc}{chapter}{1. ЛИТЕРАТУРНЫЙ ОБЗОР}

\section{Проблематика}

Многофазные течения занимают особое место в механике жидкости, поскольку в них
необходимо одновременно описывать движение нескольких сред и эволюцию границы
раздела между ними. Даже в простейшем случае двух несмешивающихся жидкостей или
жидкости с газовыми включениями задача оказывается существенно сложнее
однофазного течения: плотность и вязкость могут изменяться скачком, форма
интерфейса заранее неизвестна, а силы поверхностного натяжения локализованы в
узкой окрестности границы фаз. По этой причине задачи о каплях, пузырьках,
струях и свободных поверхностях используются как базовые тесты для развития
методов вычислительной гидродинамики.

Течения жидкостей с пузырьками возникают во многих инженерных и природных
процессах. Один из наиболее важных примеров связан с кипением и теплообменом.
При подводе тепла к жидкости образование пузырьков пара меняет локальную
теплопередачу, влияет на устойчивость температурного режима и может приводить
как к росту эффективности теплообменника, так и к опасным режимам перегрева.
В энергетике такие процессы встречаются в парогенераторах, испарителях и
системах охлаждения. Для их расчета важно знать не только среднюю долю газа, но
и форму пузырьков, критерий их отрыва от поверхности, взаимодействие с
пограничным слоем и дальнейшее всплытие.

Практическая сложность таких процессов состоит в том, что наблюдаемая
макроскопическая характеристика установки часто определяется микроструктурой
потока. Например, одинаковая средняя газовая доля может соответствовать разным
распределениям пузырьков по размерам, разной скорости подъема и разной
интенсивности перемешивания. Следовательно, для надежного прогноза
недостаточно использовать только усредненные инженерные корреляции. Требуется
модель, которая хотя бы приближенно восстанавливает локальные поля скорости,
давления и положения интерфейса. Именно поэтому даже простая задача
одиночного пузырька сохраняет методическую ценность: она позволяет изучать
связь между локальной формой интерфейса и глобальной динамикой всплытия.

Другой класс приложений связан с газообразованием и барботажем. В химических
реакторах газ вводится в жидкую фазу для перемешивания, насыщения реагентом или
интенсификации массообмена. Эффективность такого процесса определяется
размером пузырьков, скоростью их подъема, временем пребывания газа в жидкости и
площадью интерфейса. Если пузырьки слишком быстро объединяются,
площадь интерфейса уменьшается; если они дробятся, возрастает сопротивление и
усложняется структура потока. Поэтому модели пузырьковой динамики используются
при проектировании колонных аппаратов, ферментеров, флотационных установок и
других систем, где газовая и жидкая фазы должны интенсивно взаимодействовать.

Для промышленных аппаратов важна не только возможность один раз рассчитать
режим, но и способность быстро оценивать последствия изменения параметров:
расхода газа, вязкости жидкости, размера начальных включений или интенсивности
перемешивания. В режиме проектирования это приводит к большому числу
параметрических расчетов. В режиме эксплуатации возникает близкая задача: по
ограниченным измерениям необходимо понять, как изменяется внутреннее состояние
аппарата и не приближается ли он к нежелательному режиму. Прямое измерение
полей скорости и давления внутри непрозрачной многофазной среды обычно
затруднено, поэтому возрастает роль расчетных и гибридных методов, которые
могут восстанавливать недоступные величины по косвенным данным.

Близкие задачи возникают при описании эмульсий и дисперсных систем. В эмульсии
одна жидкость распределена в другой в виде капель, и ее свойства зависят от
размерного распределения, деформации и коалесценции включений. Такие среды
важны для нефтехимии, пищевой промышленности, фармацевтики и материаловедения.
С математической точки зрения капля в жидкости и газовый пузырек имеют много
общего: в том и другом случае необходимо отслеживать движущуюся границу раздела,
учитывать поверхностное натяжение и связывать движение интерфейса с полями
скорости и давления. Отличия проявляются в соотношении плотностей, вязкостей и
возможности сжимаемости газа, однако базовые численные трудности остаются
сходными.

Задача всплытия одиночного пузырька является удобной модельной задачей для
анализа таких процессов. В ней присутствуют основные эффекты многофазной
гидродинамики: плавучесть, вязкое сопротивление, деформация интерфейса,
поверхностное натяжение и взаимодействие пузырька с индуцированным полем
скорости. При этом геометрия задачи достаточно проста, чтобы использовать ее
как воспроизводимый тест в качестве проверки численных схем. В работе Hysing и соавторов
\cite{Hysing2009} двумерная задача о всплытии пузырька была предложена как
количественный тест для сравнения разных CFD-методов. Такой тест позволяет
оценивать не только форму интерфейса, но и траекторию центра масс, скорость
подъема и чувствительность решения к дискретизации.

Основная вычислительная сложность двухфазных течений связана с тем, что
наиболее важные физические процессы происходят около интерфейса. В этой области
поля быстро меняются, а малое смещение границы фаз может привести к заметной
ошибке в распределении плотности, давления и скорости. Для прямого численного
моделирования обычно требуется достаточно мелкая сетка, особенно если требуется
учитывать кривизну интерфейса и поверхностное натяжение. При увеличении
размерности, времени расчета или числа параметров задачи вычислительная сложность CFD-решений
быстро возрастает.

В последние годы для таких задач стали рассматривать модели машинного обучения,
которые не заменяют физические законы полностью, а используют их как часть
обучения. Физико-информированные нейронные сети (Physics-Informed Neural
Networks, PINN) вводят невязки, описывающие дифференциальные уравнения, а также начальные и граничные условия, в функцию потерь и,
тем самым, связывают аппроксимацию с законами сохранения \cite{Raissi2019}.
Для двухфазных течений этот подход особенно интересен в обратных задачах:
например, когда известна эволюция интерфейса, но требуется восстановить поля
скорости и давления. В работе Бухендвы, Адами и Адамса \cite{Buhendwa2021}
показано, что PINN можно применять к несжимаемым двухфазным течениям с
VOF-описанием интерфейса, включая задачу всплытия пузырька. Настоящая работа
развивает эту идею в параметрическом направлении: начальный радиус пузырька
рассматривается не как фиксированная константа, а как входной
параметр модели.

\section{Основные уравнения многофазной гидродинамики}

В основе большинства моделей несжимаемых двухфазных течений лежат уравнения
сохранения массы и импульса. Условие несжимаемости жидкости записывается в виде
\begin{equation}
    \nabla\cdot\mathbf{u}=0.
    \label{eq:review_incompressibility}
\end{equation}
Движение жидкости описывается уравнением Навье--Стокса:
\begin{equation}
    \rho
    \left(
        \frac{\partial\mathbf{u}}{\partial t}
        +(\mathbf{u}\cdot\nabla)\mathbf{u}
    \right)
    =
    -\nabla p
    +\nabla\cdot
    \left[
        \mu
        \left(\nabla\mathbf{u}+\nabla\mathbf{u}^{T}\right)
    \right]
    +\rho\mathbf{g}
    +\mathbf{f}_{\sigma}.
    \label{eq:review_momentum}
\end{equation}
Здесь \(p\) обозначает давление, \(\mathbf{g}\) --- ускорение свободного
падения, а \(\mathbf{f}_{\sigma}\) --- силу поверхностного натяжения. В
многофазной задаче коэффициенты \(\rho\) и \(\mu\) зависят от того, какая фаза
находится в данной точке пространства.

Характер движения пузырька определяется балансом нескольких физических
механизмов. Сила Архимеда задает направление всплытия, вязкие напряжения
ограничивают скорость движения, а поверхностное натяжение стремится уменьшить
площадь интерфейса и сохранить более компактную форму. Если вязкость велика,
пузырек поднимается медленно и деформируется слабо. При малом вязком
сопротивлении и большом влиянии инерции форма может заметно вытягиваться, а в
следе за пузырьком формируется более сложное вихревое течение. Поэтому одна и
та же математическая система может описывать разные режимы, а качество
численного метода следует проверять не только по полю одной выбранной величины, но и по
совокупности признаков: форме интерфейса, скорости подъема, распределению
давления и структуре скорости.

Один из распространенных способов описания динамики фаз состоит во введении
индикаторной или объемной переменной. В методе Volume of Fluid (VOF) используется
объемная доля \(\alpha\): значение \(\alpha=1\) соответствует одной фазе,
\(\alpha=0\) --- другой фазе, а промежуточные значения возникают в ячейках,
через которые проходит интерфейс \cite{Hirt1981}. Тогда материальные параметры
можно записать как сглаженную или кусочно-линейную смесь:
\begin{align}
    \rho(\alpha) &= \alpha\rho_1+(1-\alpha)\rho_2, \\
    \mu(\alpha) &= \alpha\mu_1+(1-\alpha)\mu_2.
\end{align}
Если фазовые переходы и растворение не рассматриваются, объемная доля
переносится потоком:
\begin{equation}
    \frac{\partial\alpha}{\partial t}
    +\mathbf{u}\cdot\nabla\alpha=0.
    \label{eq:review_alpha_transport}
\end{equation}
Это уравнение связывает движение интерфейса с полем скорости. Оно также
показывает одну из главных сложностей задачи: численная схема должна переносить
резкий профиль \(\alpha\) без чрезмерного размытия и без возникновения
нефизических осцилляций.

В реальном расчете поле \(\alpha\) выполняет сразу две функции. С одной
стороны, оно показывает положение фаз и тем самым задает геометрию интерфейса.
С другой стороны, через него вычисляются плотность, вязкость и поверхностные
силы. Ошибка в переносе \(\alpha\) поэтому не остается локальной ошибкой
визуализации: она меняет коэффициенты уравнения импульса и влияет на дальнейшую
эволюцию всего течения. Это особенно важно в задачах с большим отношением
плотностей или вязкостей, где небольшое смещение границы фаз приводит к
заметному изменению локальной инерции и диссипации.

Поверхностное натяжение обычно связывают с кривизной интерфейса. В
континуальной формулировке CSF (Continuum Surface Force) сила поверхностного
натяжения распределяется по узкой окрестности интерфейса и записывается через
кривизну \(\kappa\) и нормаль к границе фаз \cite{Brackbill1992}:
\begin{equation}
    \mathbf{f}_{\sigma}
    =
    \sigma\kappa\nabla\alpha,
    \qquad
    \kappa=-\nabla\cdot
    \frac{\nabla\alpha}{|\nabla\alpha|}.
    \label{eq:review_csf}
\end{equation}
Коэффициент \(\sigma\) задает поверхностное натяжение. Такая запись удобна для
методов, где интерфейс представлен полем на сетке, однако она требует
аккуратного вычисления градиентов и кривизны. Ошибка в кривизне может вызывать
паразитные токи около интерфейса, что особенно заметно в задачах с большим
вкладом поверхностного натяжения \cite{Popinet2018}.

В литературе существуют и другие способы описания многофазности. В
методах уровней интерфейс задается нулевым уровнем вспомогательной функции,
что удобно для вычисления нормалей и кривизны \cite{Osher1988}. В
методах отслеживания перемещения фронта граница раздела хранится явно в виде лагранжевой сетки,
движущейся вместе с потоком \cite{Tryggvason2001}. В методах с диффузной границей и
моделях фазовых полей переход между фазами намеренно делается гладким, а
структура интерфейса описывается дополнительной энергией смешения
\cite{Anderson1998}. Выбор конкретной формулировки зависит от того, что важнее
в рассматриваемой задаче: строгая консервативность массы, точная геометрия
интерфейса, устойчивость при топологических изменениях или удобство реализации
в численном методе.

\section{Методы численной гидродинамики (CFD)}

Классические CFD-подходы к многофазным течениям можно рассматривать на
нескольких уровнях. Первый уровень связан с дискретизацией уравнений сохранения
на сетке: сюда относятся методы конечных разностей, конечных объемов и конечных
элементов. Отдельный класс составляют бессеточные лагранжевы методы, например, метод гидродинамики гладких частиц (SPH), где среда представляется набором движущихся частиц. Второй уровень связан
с тем, как в расчете представляется интерфейс: через объемную долю,
функцию уровня, явную поверхность или диффузный интерфейс. Оба выбора важны
одновременно. Например, VOF-метод задает способ хранения интерфейса, но на
практике часто реализуется внутри конечно-объемной схемы для уравнений
Навье--Стокса.

Сравнивая перечисленные методы, нельзя сводить их вычислительную стоимость
лишь к числу ячеек сетки. В нестационарной двухфазной постановке приходится
брать маленький шаг по времени, совместно решать уравнения для давления и
скорости, переносить фазовую переменную, считать поверхностное натяжение и
следить за устойчивостью схемы. Когда интерфейс перемещается с большой
скоростью или заметно деформируется, геометрию приходится пересчитывать на
каждом шаге. Из-за этого затратной оказывается уже серия расчетов плоской
модельной задачи, а для трехмерных инженерных постановок, как правило, нужны
специализированные вычислительные мощности.

\textbf{Метод конечных разностей (FDM).}
В основе метода конечных разностей лежит замена производных из
дифференциальных уравнений их разностными аналогами, выраженными через
значения искомых величин в соседних узлах. Так, производные скорости, давления
или фазовой переменной записываются разностными формулами нужного порядка
точности. Реализуется такой подход сравнительно легко и естественно ложится на
регулярные прямоугольные сетки; объем вычислений при этом растет
пропорционально числу узлов и числу шагов по времени, а сами разностные
операторы имеют удобную для расчета структуру. Слабое место метода --- сильная
привязка к геометрии сетки: в областях сложной формы и у искривленных границ
требуются либо преобразования координат, либо специальные шаблоны. В
многофазных течениях этот недостаток проявляется особенно остро: резкий
интерфейс и скачки коэффициентов легко порождают численные осцилляции, если
схема недостаточно устойчива или плохо согласована с переносом фазовой
переменной.

\textbf{Метод конечных объемов (FVM).}
Метод конечных объемов опирается на интегральную запись законов сохранения.
Расчетную область делят на контрольные объемы, и изменение массы, импульса или
иной величины внутри каждой ячейки задается потоками через ее грани. Ключевое
достоинство здесь --- локальная консервативность: для всякой ячейки баланс
сохраняемой величины выдерживается с точностью до потоков на ее границах.
Именно поэтому метод так распространен в вычислительной гидродинамике и
естественно стыкуется с VOF-подходом, где принципиально важно сохранять объем
каждой фазы. Вычислительная сложность метода зависит
от числа контрольных объемов, порядка аппроксимации потоков и стоимости решения
системы для давления и скорости. Метод может работать как на структурированных,
так и на неструктурированных сетках, но качество решения все равно зависит от
разрешения около интерфейса и от численной диффузии в схемах переноса. В
задачах с поверхностным натяжением конечно-объемная схема должна также
аккуратно согласовывать давление, кривизну и фазовую переменную, иначе около
интерфейса могут появляться паразитные скорости.

\textbf{Метод гладких частиц (SPH).}
Метод гладких частиц относится к бессеточным лагранжевым
методам. В этом подходе жидкость представляется набором частиц, которые
переносят массу, скорость, плотность и другие величины. Значения полей и их
производные восстанавливаются с помощью сглаживающего ядра по соседним
частицам. Такой способ описания удобен для задач с сильными деформациями
свободной поверхности, разрывом струй, перемешиванием и движением сложных
границ, поскольку расчетная сетка не требуется и не возникает задача ее
перестроения. Для многофазных течений SPH привлекателен тем, что интерфейс
между фазами естественно перемещается вместе с частицами.

Недостатки SPH связаны с точностью и вычислительной сложностью. Для каждой
частицы необходимо определять соседей в радиусе сглаживания, поэтому вычислительная сложность прямого алгоритма весьма высока, а эффективная реализация требует специальных структур поиска. При
увеличении числа частиц стоимость расчета растет быстро, особенно в трехмерных
задачах. Точность метода зависит от распределения частиц, выбора ядра,
длины сглаживания и стабилизирующих слагаемых. В задачах
гидродинамики несжимаемой жидкости требуется отдельно контролировать плотность и давление, иначе
могут возникать шум давления, неустойчивости распределения частиц и
нефизические колебания около интерфейса. Поэтому SPH хорошо подходит для
сильно деформируемых свободных поверхностей, но при расчете тонких эффектов
поверхностного натяжения и точных полей давления часто требует аккуратной
калибровки и большого числа частиц.

После выбора базовой дискретизации необходимо выбрать способ представления
границы раздела фаз. Именно этот выбор определяет, как в расчете задаются
положение пузырька, кривизна интерфейса, поверхностное натяжение и перенос
фазовой переменной.

\textbf{Метод Volume of Fluid (VOF).}
Реализация VOF-метода предусматривает хранение в каждой ячейке сетки объемной доли одной из фаз
\cite{Hirt1981}. Его главное преимущество состоит в консервативности: при
корректной дискретизации суммарный объем фазы сохраняется достаточно хорошо
даже при сильной деформации интерфейса. Поэтому VOF часто используют в задачах
с каплями, пузырьками, струями и свободными поверхностями. Вычислительная
сложность метода в основном определяется решением уравнений Навье--Стокса и
переносом \(\alpha\); дополнительные затраты связаны с реконструкцией
интерфейса и вычислением кривизны. Метод сильно зависит от качества сетки около
границы фаз: если сетка грубая, интерфейс размывается, а кривизна становится
неточной. Устойчивость VOF-подхода обычно хорошая при соблюдении ограничений на
шаг по времени, но поверхностное натяжение может порождать паразитные скорости,
если геометрия интерфейса восстановлена недостаточно точно.

Важным достоинством VOF является возможность работать с разрывом фазовой
переменной без явного хранения отдельной поверхности. Это удобно при слиянии и
разделении пузырьков. Однако за это приходится платить сложностью
геометрической реконструкции. Для получения качественной формы интерфейса часто
используют специальные схемы, которые восстанавливают положение границы внутри
ячейки. Если такая реконструкция недостаточно точна, ошибка проявляется не
только в поле \(\alpha\), но и в силе поверхностного натяжения, поскольку она
зависит от кривизны.

\textbf{Метод уровней (Level-set).}
В методе уровней интерфейс отождествляют с нулевым множеством вспомогательной
функции \(\phi\) \cite{Osher1988}. Как правило, \(\phi\) задают близкой к
знаковой функции расстояния, и тогда нормаль и кривизна границы получаются
прямым дифференцированием. Геометрию произвольной формы такой способ описывает
удобно, а слияние и распад интерфейса обрабатываются сами собой. По объему
вычислений метод близок к VOF, однако дополнительно нужно интегрировать
уравнение переноса для \(\phi\) и время от времени проводить реинициализацию,
иначе свойства функции расстояния теряются. Основная слабость в том, что в
исходной формулировке метод не консервативен строго: на длинных интервалах
объем фазы способен медленно дрейфовать. К сетке метод чувствителен через
точность вычисления градиентов и кривизны. Устойчивость, как правило, хорошая,
но компромисс между частотой реинициализации и сохранением массы приходится
подбирать аккуратно.

В силу этого метод уровней принято считать геометрически выигрышным, но
нуждающимся в отдельных средствах контроля объема. На практике его часто
связывают с VOF: уровневая функция отвечает за гладкую форму интерфейса, а
VOF-переменная удерживает массу фазы. Гибридные схемы такого рода поднимают
качество решения ценой усложнения алгоритма и роста числа подгоночных
параметров.

\textbf{Метод отслеживания фронта.}
Здесь граница фаз задается напрямую: ее описывают цепочкой лагранжевых
маркеров или отдельной сеткой, которая движется вместе с потоком
\cite{Tryggvason2001}. За счет этого положение интерфейса известно с высокой
точностью, а поверхностное натяжение считается прямо по геометрии границы.
Метод выигрышен там, где топология интерфейса меняется нечасто, а к форме
предъявляются жесткие требования. Платой служит более высокая вычислительная
стоимость по сравнению с чисто эйлеровыми схемами: приходится постоянно
согласовывать лагранжев интерфейс с эйлеровой сеткой, на которой считаются
скорость и давление. Зависимость от разрешения двоякая --- итог определяется и
основной сеткой течения, и густотой маркеров на интерфейсе. Пока граница
эволюционирует плавно, метод устойчив, но слияние и разрыв пузырьков требуют
отдельных процедур перестройки топологии.

Когда пузырек один, явное отслеживание интерфейса оказывается весьма выгодным:
поверхность остается связной, и переносить ее вместе с потоком несложно. А вот
в потоках с множеством включений сопровождение интерфейсной сетки превращается
в трудоемкую задачу --- нужно ловить сближение поверхностей, не давать
элементам вырождаться и решать, когда слияние или дробление физически
оправдано. Сделать эти операции по-настоящему универсальными непросто, поэтому
метод фронта обычно берут там, где принципиальна точность геометрии и структура
интерфейса заранее ясна.

\textbf{Метод фазовых полей и подход диффузной границы.}
В фазово-полевых моделях граница между фазами перестает быть математически
резким скачком --- ее место занимает тонкий, но имеющий конечную ширину
переходный слой \cite{Anderson1998}. Этот слой описывают параметром порядка,
эволюцию которого обычно задают уравнениями типа Кана--Хилларда или
Аллена--Кана. Сильные стороны подхода --- хорошая гладкость постановки и
естественная обработка перестроек топологии. Вдобавок поверхностное натяжение
вытекает из энергетической формулировки, что упрощает анализ устойчивости.
Обратная сторона --- необходимость разрешать толщину диффузного интерфейса:
слишком тонкий слой вынуждает измельчать сетку, а слишком толстый уводит модель
от задачи с резкой границей. Поэтому расчет может обходиться дорого, особенно в
трех измерениях и при сильном контрасте свойств фаз.

В плане устойчивости диффузная граница удобна тем, что снимает разрыв
коэффициентов на уровне самой математической модели. Однако с физическим
смыслом толщины переходного слоя надо обращаться осторожно: выбранный из
соображений счета слой может оказаться шире реальной зоны интерфейса. Решение
при этом сглаживается, но заодно теряется часть эффектов резкой границы. Для
инженерных приложений отсюда следует, что чувствительность результата к
параметрам интерфейсной модели нужно специально проверять.

\textbf{Адаптивные сетки и специализированные CFD-решатели.}
В пузырьковых задачах нередко прибегают к адаптивному сгущению сетки вблизи
интерфейса. Смысл в том, чтобы не наращивать число ячеек по всей области, а
выделять мелкое разрешение лишь там, где велики градиенты \(\alpha\), скорости
или давления. По сравнению с равномерно мелкой сеткой это экономит ресурсы, но
усложняет код: сетку надо перестраивать, переносить данные между уровнями и
следить за сохранением физических величин. Для задач с поверхностным
натяжением адаптация особенно уместна, ведь ошибка кривизны сосредоточена возле
границы фаз. Тем не менее устойчивость по-прежнему упирается в шаг по времени,
качество реконструкции интерфейса и согласование давления, вязкости и
капиллярных сил.

Получать эталонные данные адаптивными методами удобно: они дают возможность
резко поднять разрешение в критических зонах. В рамках настоящей работы такие
данные служат опорой для обучения и проверки PINN-модели. Но адаптивность сама
по себе расчет не ускоряет: стоит изменить параметры задачи, и моделирование
обычно запускают с нуля, а время до ответа все так же определяется числом шагов
по времени и сложностью решателя.

В итоге классические CFD-методы дают подробную и физически осмысленную картину
двухфазного течения, но при многократном использовании обходятся
вычислительно дорого. Особенно остро это проявляется в параметрических
исследованиях, когда моделирование приходится повторять для каждого отдельного
набора параметров --- при разных радиусах пузырька, числах Рейнольдса,
коэффициентах поверхностного натяжения или свойствах фаз. На каждый новый набор
обычно уходит отдельный нестационарный расчет с большим числом шагов по
времени. Поэтому классические методы хороши как средство глубокого анализа и
источник качественных данных, но мало пригодны для быстрого перебора режимов и
для прямого встраивания в установки, работающие в реальном времени.

Для задач управления и диагностики это ограничение особенно чувствительно. На
действующей установке оператор или автоматика должны оценивать состояние
быстрее, чем заметно меняется сам процесс. Когда от модели ждут помощи в выборе
расхода газа, предупреждения о сваливании в нежелательный режим, оценки
положения интерфейса или восстановления поля скорости по неполным наблюдениям,
задержка в минуты или часы делает прямой CFD-расчет непригодным. К тому же в
промышленной системе зачастую недопустимо гонять тяжелое моделирование на
каждое новое измерение или малейшее изменение параметров.

Отсюда вытекает интерес к приближенным моделям, дополняющим CFD: они опираются
на результаты численных расчетов и на ограничения, налагаемые уравнениями
гидродинамики. PINN-подход относится именно к такому классу:
он не отменяет роль базовых CFD-данных, но позволяет включить физические
невязки в обучение и получить непрерывную аппроксимацию полей в
пространстве-времени. После обучения нейронная сеть может выдавать приближенное
решение существенно быстрее, чем полный численный расчет. Если ввести в сеть
параметры задачи, например начальный радиус пузырька, такая модель может
использоваться как быстрая суррогатная модель для семейства близких режимов. В
этом состоит прикладная мотивация настоящей работы: параметрическая PINN может
быть рассмотрена как элемент цифрового двойника или расчетного модуля для
быстрой оценки динамики пузырька, где классический CFD остается источником
обучающих данных и эталоном точности, но не используется непосредственно на
каждом шаге работы установки.

\section{Физико-информированные нейронные сети}

Физико-информированные нейронные сети (Physics-Informed Neural Networks --- PINN)
были предложены Рейсси, Пердикарисом и Карниадакисом \cite{Raissi2019} как
универсальный подход для решения прямых и обратных задач, описываемых
дифференциальными уравнениями в частных производных. Ключевая идея состоит в
том, что нейронная сеть обучается не только воспроизводить данные наблюдений,
но и удовлетворять физическим законам, которые вводятся в функцию потерь в виде
невязок, построенных на основе дифференциальных уравнений. Это позволяет совместить данные измерений
(или результаты CFD-расчётов) с ограничениями, накладываемыми физикой задачи.

Формально пусть искомые поля $\mathbf{u}(\mathbf{x}, t)$ описываются системой
дифференциальных операторов $\mathcal{N}[\mathbf{u}]=0$ на некоторой области
$\Omega$ с граничными условиями $\mathcal{B}[\mathbf{u}]=0$ на $\partial\Omega$.
Нейронная сеть $\hat{\mathbf{u}}_\theta(\mathbf{x}, t)$ с параметрами $\theta$
аппроксимирует искомые поля. Функция потерь при обучении складывается из четырёх
слагаемых:
\begin{equation}
\mathcal{L}(\theta) = \lambda_{\mathrm{data}}\mathcal{L}_{\mathrm{data}}(\theta)
  + \lambda_{\mathrm{pde}}\,\mathcal{L}_{\mathrm{pde}}(\theta)
  + \lambda_{\mathrm{bc}}\,\mathcal{L}_{\mathrm{bc}}(\theta)
  + \lambda_{\mathrm{ic}}\,\mathcal{L}_{\mathrm{ic}}(\theta),
\end{equation}
где $\mathcal{L}_{\mathrm{data}}$ --- среднеквадратичное отклонение предсказаний
от обучающих данных, $\mathcal{L}_{\mathrm{pde}}$ --- средняя невязка уравнений
на коллокационных точках внутри области, $\mathcal{L}_{\mathrm{bc}}$ и $\mathcal{L}_{\mathrm{ic}}$ ---
невязка граничных и начальных условий соответственно, а $\lambda_{*}$
--- веса соответствующих слагаемых. Производные для вычисления невязок PDE
находятся через автоматическое дифференцирование графа вычислений нейронной сети.

Среди ключевых преимуществ подхода следует выделить несколько. Во-первых,
PINN не привязан к фиксированной сетке: коллокационные точки могут произвольно
располагаться в пространстве-времени, что делает метод особенно удобным для
задач с подвижными границами. Во-вторых, встраивание физических законов в
функцию потерь играет роль регуляризатора: даже при ограниченном числе
обучающих данных сеть вынуждена формировать физически согласованные поля.
В-третьих, PINN естественно справляется с обратными задачами, когда часть
физических параметров неизвестна и определяется совместно с полями в процессе
обучения. В-четвёртых, после обучения вывод нейронной сети занимает доли
секунды, что на несколько порядков быстрее полного CFD-расчёта.

Вместе с тем метод имеет существенные ограничения. Стоимость обучения высока:
требуется многократное выполнение прямого и обратного распространения ошибки с вычислением
производных высокого порядка, что делает одно обучение сравнимым по времени
с серией CFD-симуляций. Сходимость чувствительна к выбору весов слагаемых в функции потерь и к
распределению коллокационных точек. При резких скачках параметров, как это
происходит на интерфейсе в двухфазных задачах, нейронные сети с
ограниченным числом слоёв и малым количеством нейронов в скрытых слоях могут испытывать затруднения с
аппроксимацией.

Применение PINN к несжимаемым двухфазным течениям, в частности к задаче
всплытия пузырька, было исследовано в работе Бухендвы, Адами и Адамса
\cite{Buhendwa2021}. Авторы показали, что по наблюдаемой эволюции интерфейса,
заданной в виде поля объёмной доли (VOF), можно восстановить скрытые поля
скорости и давления. Несмотря на то, что PINN является нейронной сетью,
физические ограничения (уравнения Навье--Стокса, несжимаемость,
транспорт фазового индикатора) позволяют надёжно восстановить поля даже
при ограниченном числе граничных условий. Эта работа стала непосредственным
прототипом для настоящего исследования.

В настоящей работе подход расширяется до \emph{параметрической} PINN: радиус
пузырька $R$ передаётся в сеть как дополнительный входной признак, что
позволяет обучить одну модель сразу на нескольких режимах и затем получать
решение для любого значения начального радиуса из заданного диапазона без повторного обучения.
Такой подход сближает PINN с нейросетевыми операторами (Neural Operators),
однако сохраняет прямую интерпретацию через физические невязки в функции
потерь.
